%============================================================
% MTH 112 Project - Double Angle Identities
% Updated 202504
%============================================================

\section{Double Angle Identities} \label{section-double-angle-identities}
In this section, we will learn to use double-angle and half-angle formulas to find exact trigonometric values and verify identities, and use reduction formulas to simplify an expression.

\textbook{\href{https://openstax.org/books/algebra-and-trigonometry-2e/pages/9-3-double-angle-half-angle-and-reduction-formulas}{9.3}}

\subsection*{Preparation Exercises} \label{preparation-exercises-section-double-angle-identities}

Answer the following without using a calculator. These questions are intended to check the prerequisite skills needed to complete the rest of the lab.

\begin{myPrep}
	Use the sum of angles for sine formula, $\sin(\alpha+\beta)=\sin(\alpha)\cos(\beta)+\cos(\alpha)\sin(\beta)$, to evaluate and simplify the expression $\sin(\theta+\theta)$.
\end{myPrep}
	\vspace{3cm}

\begin{myPrep}
	In which quadrants are the function values negative?
		\begin{enumerate}
			\begin{multicols}{4}
			\item $\sin(\theta)$.
			\item $\cos(\theta)$.
			\item $\tan(\theta)$.
			\item $\sec(\theta)$.
			\end{multicols}
		\end{enumerate}
\end{myPrep}
	\vfill

\begin{myPrep}
	True or False: $\sin(60^\circ+60^\circ)=2\sin(60^\circ)$. Explain your reasoning.
\end{myPrep}
	\vfill

\begin{myPrep}
	Verify the identity $\frac{\cos(\tau)}{1+\sin(\tau)}=\sec(\tau)-\tan(\tau)$ algebraically.
\end{myPrep}
\vspace{6cm}

%======================================================
\newpage
%======================================================

\subsection*{Practice Exercises} \label{practice-exercises-section-double-angle-identities}

\begin{myPractice}
	Find the values using the half angle identities.
	\begin{enumerate}
		\begin{multicols}{3}
		\item $\sin(22.5^\circ)$
		\item $\cos(15^\circ)$
		\end{multicols}
	\end{enumerate}
\end{myPractice}
\vfill

\begin{myPractice}
	If $\sin(\phi)=\frac{1}{6}$ and $\phi$ is in quadrant II, find the following values.
	\begin{enumerate}[itemsep=3cm]
		\begin{multicols}{2}
		\item $\cos(\phi)$
		\item $\tan(\phi)$
		\item $\sin(2\phi)$
		\item $\cos(2\phi)$
		\item $\tan(2\phi)$
		\item $\sin(\frac{\phi}{2})$
		\item $\cos(\frac{\phi}{2})$
		\end{multicols}
	\end{enumerate}
\end{myPractice}
\vfill
%======================================================
\newpage
%======================================================

\begin{myPractice}
	Verify the identities algebraically.
	\begin{enumerate}
		\item $\sec^2(\theta)=\frac{2}{1+\cos(2\theta)}$
		\vfill
		\vfill
		\item $\left(2\sin(\theta)-3\cos(\theta)\right)^2=4-6\sin(2\theta)+5\cos^2(\theta)$
		\vfill
		\vfill
	\end{enumerate}
\end{myPractice}

\begin{myPractice}
	Using a graphing utility, find the value of $C$ that makes the equation $\left(\cos(x)+\sec(x)\right)^2=C+\tan^2(x)-\sin^2(x)$ an identity.
	\vfill
\end{myPractice}

%======================================================
\newpage
%======================================================

\subsection*{Definitions} \label{definitions-section-double-angle-identities}

\begin{myDefinition}[{Double-Angle Identities}]~\\[0.5mm]
	Each of the trigonometric functions has a formula that accounts for an angle that is twice as large as some starting angle, the \textbf{double-angle identities}. We will focus only on those for sine, cosine, and tangent. The formulas are summarized below.
		\begin{itemize}
			\item $\sin(2\theta)=2\sin(\theta)\cos(\theta)$
			\vfill
			\item 	$\cos(2\theta)=\cos^2(\theta)-\sin^2(\theta)$\\
					$\phantom{\cos(2\theta)}=1-2\sin^2(\theta)$\\
					$\phantom{\cos(2\theta)}=2\cos^2(\theta)-1$
			\vfill
			\item $\tan(2\theta)=\frac{2\tan(\theta)}{1-\tan^2(\theta)}$
			\vfill
		\end{itemize}
\end{myDefinition}

% FIX Issue: Include pictures to illustrate double and half angles. Make sure to cut any tan of half angles since they're nasty and maybe change the definition of tan of half angle to be the commented out version below.
\begin{myDefinition}[{Half-Angle Identities}]~\\[0.5mm]
	Each of the trigonometric functions has a formula that accounts for an angle that is half as large as some starting angle, the \textbf{half-angle identities}. We will focus only on those for sine, cosine, and tangent. The formulas are summarized below.
		\begin{itemize}
			\item $\sin\left(\frac{\theta}{2}\right)=\pm\sqrt{\frac{1-\cos(\theta)}{2}}$
			\vfill
			\item $\cos\left(\frac{\theta}{2}\right)=\pm\sqrt{\frac{1+\cos(\theta)}{2}}$
			\vfill
			\item $\tan\left(\frac{\theta}{2}\right)=\pm\sqrt{\frac{1-\cos(\theta)}{1+\cos(\theta)}}$
			% \item $\tan\left(\frac{\theta}{2}\right)=\frac{\sin(\theta)}{1+\cos(\theta)}$
			\vfill
		\end{itemize}
\end{myDefinition}

%======================================================
\newpage
%======================================================

\subsection*{Exit Exercises} \label{exit-exercises-section-double-angle-identities}


\begin{myExit}
	If $\tan(\rho)=\frac{4}{3}$ and $\rho$ is in quadrant III, find the following values.
		\begin{enumerate}
			\begin{multicols}{3}
				\item $\cos(\rho)$
				\item $\sin(2\rho)$
				\item $\sin(\frac{\rho}{2})$
			\end{multicols}
		\end{enumerate}
\end{myExit}
\vfill

\begin{myExit}
	Algebraically verify the identity $\sin(4\alpha)=8\sin(\alpha)\cos^3(\alpha)-4\sin(\alpha)\cos(\alpha)$.
\end{myExit}
\vfill
\exitlikert{double angle identities}